\section{最终科研判断与可发表贡献的选择}
\subsection{是否值得成为独立 IEEE TAP 主线}
\textbf{有条件值得，但当前应立项为范围受限的理论与机制研究。} 未知环境给目标推断带来的传播歧义，是独立于材料 MAE、局部 current 近似和单点重建的科学对象。这个对象与天线采集、全波多次散射、有限孔径和实际穿透损耗直接相关，具备 TAP 的物理基础。但“通用未知墙体自主递进成像框架”目前缺少三个必要证据：有限尺寸复杂墙中的重要传播歧义、可实际覆盖这些解释的推断方法、以及同费用目标风险改善。

因此本报告建议先给路线 I 一个有硬上限的 P0，而不是直接给整套算法一个长期项目。若 P0 成功，再把一个最有用的风险界与最小概率实验合成论文；若只有理想单频板的歧义，应发表或记录有限结论，而不是用复杂系统包装它。

\subsection{六个候选命题的独立性排序}
\begin{longtable}{p{23mm}p{64mm}p{61mm}}
\toprule
候选 & 本轮能够建立的内容 & 仍决定原创性与价值的缺口\\
\midrule
I：下游可辨识性 & 纤维判据、局部必要/条件充分性、非线性歧义模数、严格一维见证及二点风险下界。 & 在普通有限支撑复杂墙、有限单侧宽带和 nuisance 下给非平凡构造或上界。最高优先。\\
II：传播等价类 & 任务伪度量、四块签名、联合 prior 不变时的 likelihood--posterior 压缩界。 & 找到明显少于材料模式数的传播代表，同时保留目标条件依赖；证书必须实际可计算。\\
III：安全更新 & 任务 Bayes risk、尾质量与错误固定机制；明确确定性模型不能逼近连续材料全后验。 & 非保守的在线证书和模式覆盖。未覆盖质量未知时，不能称绝对安全。\\
IV：A17 任务反馈 & 可复用双侧增广公式、paired support、局部更新四块端口分解。 & 胜过普通 goal-ROM、同空间 Galerkin、优化 direct/Krylov 的实际选空间价值；已有 A17-X/F 覆盖大量代数。\\
V：自适应前沿 & 随机 Schur 四元组、可回退可信任务集合；不依赖预设平面层。 & 计算定义不自动形成新物理定理；需证预算调度对同一 posterior 有独立增量。\\
VI：传播驱动搜索 & 用任务响应新颖度代替像素新颖度，约束完整数据和合法 prior。 & 极端可行点不等于重要概率模式；全局拓扑发现与质量估计仍未解决。\\
\bottomrule
\end{longtable}

最有希望的组合是 \textbf{I 为核心，II 为建设性结论}：先刻画观测不能消除哪些传播差异，再证明哪些差异允许被压缩。III 是其决策解释，VI 是寻找反例和模式的实验工具。IV、V 不必进入首篇论文标题。

\subsection{三条主路线的投入与收益判断}
\paragraph{路线 I：可辨识性与不可能性论文。}
理论对象清楚，能以少量严格见证和数学界检验，而不依赖大规模 posterior sampler。最好的结果不是再给一个反射/透射不同的例子，而是证明一个随频率、孔径、噪声和 nuisance 变化的任务稳定性边界，且构造满足普通材料和全波约束的有限扰动。数值门槛最低，原创潜力最高；失败风险是物理先验和完整联合数据消除所有实用尺度歧义，或所得现象完全落在已有多模逆散射/Marchenko 条件讨论之内。此时仍可得到采集充分性结论，但须改变主张。

\paragraph{路线 II：概率 Green 更新与递进算法论文。}
最终可能最接近用户希望的成像系统，但它需要可靠模式质量、目标条件分布和复杂拓扑支持。其最强卖点应是相同观测与费用下目标风险或概率准确度改善；仅有更多候选、更小数据残差或更漂亮图像不足。当前 P2/F3 证据说明全局搜索与在线选择是瓶颈，直接上 Flow/Diffusion 不能解决这个逻辑缺口。先在可独立积分的低维分支基准上验证概率压缩，再扩展全材料场，是更可信的门槛。若多假设有效而推进无效，应把论文改为不确定环境的 joint inference，不保留无收益模块。

\paragraph{路线 III：A17 任务导向计算电磁论文。}
可直接复用最多现有代码与理论，但独立性审计最严格：A17-X 已有有限反馈端口形式，Research F 已有对偶误差接口，标准 goal-ROM 已优化输出而非全场。新的价值必须在具体选空间原则、可计算误差证书、或经总费用验证的 Pareto 优势。可先用一个目标域、几个合法传播分支和固定 rank 做微基准。若普通 Galerkin 同空间完全一致，而 A17 选空间更贵，必须停止“反馈代数自身更高效”的叙事。小核很精确也不能证明墙体正确。

\paragraph{可选路线 IV：概率 redatuming。}
将无法由有限单侧数据确定的初始聚焦/前向传播作为随机物理状态，研究内部 Green 响应的可辨部分，可能连接到成熟 Marchenko 文献。其优势是问题与波动物理贴近；风险是如果最终仍需不可获得的透射或已知 direct field，就离开原问题。建议作为路线 I 的文献和应用分支，不在当前并行启动第二套大算法。

以上新颖性和 TAP 适配性是本轮判断，不是期刊接受保证。真正论文必须收紧材料、采集与任务三类假设，并把结论限制到实际证明和实验覆盖的范围。

\subsection{与现有 A17、B1/C1 的独立科学关系}
A17 主要回答在固定材料状态和特定局部目标下，如何改进 current approximation 及材料步；B1/C1 的核心是材料多解表示、搜索和概率建模。Research W 应回答一个不同的问题：\emph{这些材料解释对指定下游电磁任务是否可互换，以及何时环境不确定性必然转化成目标风险。}

这个对象足够独立，但并非与既有研究毫无重叠。端口传播代数、finite resolvent、强方向积分和 posterior 误差接口已经有内部先例；应该明确复用和引用。独立性应来自 Maxwell 任务可辨识性的新物理结果，或传播等价类压缩的实际可计算性，不能来自把同一个有限反馈公式换变量。

已有负结果也提供了清楚边界：A30 的 oracle 表示 headroom 不能变成在线选择能力；P2 的有限候选不能变成全域 posterior；F3 的局部 continuation 不能变成跨拓扑采样；Research B 的新数据收益不能归因于固定数据递进；A26 的反馈残差改善不能变成材料/任务效果保证。A17 的最新有限非线性正例应保留，但不抵消这些限制。

\subsection{只能选一条路线时，推荐先证明什么}
推荐把首个目标写成以下有限而可证伪的命题：
\begin{quote}
给定非磁、被动、有限支撑的二维 TM 材料类，固定的单侧多静态频率集合、噪声与 nuisance prior，构造两个完整合法世界，使联合数据白化距离不超过 1，而一个预先声明的目标任务差异达到物理阈值；同时给出这一差异随带宽/孔径变化的稳定范围或消失条件。
\end{quote}
这是存在性与条件敏感性任务，不要求首轮估计其真实 posterior 质量。若成功，再问该类世界是否有足够 prior/posterior 质量，从而值得概率算法处理。若只有 calibration 歧义而 joint 数据能分开，应转而研究“目标作为环境探针”及联合采集如何改善传播可辨识性，不把它当失败数据删掉。

数学上优先寻找 $\omega_T(\delta)$ 的可计算下界与有限非线性实现，而不是又做一次局部小奇异值扫描。建设性后续则是把四块签名误差、共振裕度与尾质量合成非空的后验压缩容差。若上界总是大于 1 或求界比 full solve 更贵，它仍是正确数学，却没有算法证书价值。

\subsection{下一轮最小实验与停止规则}
\textbf{下一轮只执行 P0-B：} 一个公开固定的倾斜厚度变化介质体加不规则空气区、全域材料表示、同侧 4 发 12 收、两种预定带宽、两个固定背景种子。先寻找 calibration 近等价/传播不同见证，随后强制通过 target-present 联合检验、独立网格和 nuisance 合法性。普通 adjoint/Full-J 与传播危险方向提案享有相同求解器、正确器和完整预算。所有失败状态和最好找到的 pair 都保留。

出现以下结果时，应停止或收缩对应主线，而非继续加网络和搜索规模：
\begin{enumerate}
\item 传播歧义仅存在于无限平面单频精细调谐，在预定有限宽带 joint 实验中没有实用见证：停止通用未知墙递进叙事，保留边界结果。
\item 合法近数据世界对目标任务基本等效：材料多解不是该任务主要瓶颈，停止为传播分支增加复杂度。
\item 明显传播歧义存在，但仅靠 truth-selected candidate/weights 才改善：算法 NO-GO，先报告可辨识性和搜索限制。
\item 压缩证书无法避免未知真值、缺失模式质量或高成本全参考：不称在线安全更新，保留 offline audit 身份。
\item 同预算直接 joint inference 已达到相同或更好目标风险：停止递进增量主张；如多解仍有价值，保留多解方法。
\item A17 在强 direct/Krylov/goal-ROM 对照下无总费用优势：停止此计算分支，不否认其已有局部机制。
\item Oracle 真墙下目标也不可测：当前采集不支持该目标任务。新增传感/频率是另一个明确采集问题。
\end{enumerate}

本轮已经足以否定“只要逐步恢复材料并更新 Green，就会自然可靠地看得更深”的无条件主张；尚不足以否定所有有结构先验和适当采集的递进方法。值得继续检验的是：\emph{对哪些任务，环境解释仍有实质传播分歧；对哪些任务，这些分歧已经小到可以被认证地合并。} 这比承诺普适逐层恢复更窄，也更接近可验证的独立科学贡献。
